\section{问题三：偏差检验驱动的稳健融合}

问题三首先估计时间偏差，再判断两类实际观测之间是否存在需要显式建模的系统偏差。为处理时间相关残差，本文不采用独立同分布假设下的普通均值检验，而使用 HAC 协方差构造二维 Wald 统计量，并辅以移动块 bootstrap 置信区间和 HAC 趋势检验。

设对齐后的两源位置差为 $\bm d_t$，检验
\[
H_0:\ E(\bm d_t)=\bm 0.
\]
在显著性判断之外，引入偏差模相对于稳健噪声尺度的工程效应指数。只有当统计显著且工程效应达到预设阈值时，才启用系统偏差状态；若同时检出显著趋势，再允许偏差随机游走。

正式数据下统一时间偏差为 $\delta=\QThreeDelta\,\mathrm{s}$，HAC--Wald $p=\QThreeWaldP$，偏差模约 \QThreeBiasNorm\,m，工程效应指数为 \QThreeEffectIndex。移动块 bootstrap 的两个坐标区间均覆盖 0，趋势检验亦未发现显著漂移，因此正式融合不启用系统偏差状态。

随后在校正后的原始异步观测事件上运行稳健 Kalman 滤波和 RTS 平滑，并传播状态协方差得到 10 Hz 位置、速度、加速度及模型内位置标准差。该结论应表述为“当前样本与阈值下未发现具有统计和工程意义的相对系统偏差”，而不是证明真实偏差严格等于零。

\input{generated/table_q3.tex}
